\documentclass[10pt]{amsart}

\usepackage[margin=0.95in]{geometry}
\usepackage{amsmath,amssymb,mathtools}
\usepackage{booktabs}
\usepackage{microtype}
\usepackage[hidelinks]{hyperref}

\hypersetup{
  pdftitle={Counterexamples to Conjecture 5.4 of Parisi--Spahiu--Skandera--Wang},
  pdfauthor={Owen Shen},
  pdfkeywords={Kazhdan--Lusztig basis, Hecke algebra, reversal factorization, pattern avoidance}
}

\newtheorem{theorem}{Theorem}
\newtheorem{corollary}[theorem]{Corollary}
\theoremstyle{remark}
\newtheorem{remark}[theorem]{Remark}

\newcommand{\Sn}{\mathfrak S_n}
\newcommand{\Ssix}{\mathfrak S_6}
\newcommand{\Ct}{\widetilde C}
\newcommand{\Cab}[2]{C_{#1#2}}
\newcommand{\gap}{\operatorname{gap}_{3412}}

\title[Counterexamples to Conjecture 5.4]
{Counterexamples to Conjecture 5.4 of
Parisi--Spahiu--Skandera--Wang}

\author{Owen Shen}
\address{Operations Research Center, Massachusetts Institute of Technology,
Cambridge, Massachusetts, USA}
\email{owenshen@mit.edu}
\date{}

\subjclass[2020]{Primary 05E10; Secondary 20C08}
\keywords{Kazhdan--Lusztig basis, Hecke algebra, reversal factorization,
pattern avoidance}

\begin{document}

\begin{abstract}
Parisi, Spahiu, Skandera, and Wang conjectured that if a permutation
$w$ contains the pattern $45312$, then the modified signless
Kazhdan--Lusztig basis element $\Ct_w$ has no reversal factorization.
We give exact identities in $H_6(q)$ showing that seven such
permutations admit reversal factorizations, and we prove that these are
exactly the factorable $45312$-containing elements of $\Ssix$.
The other nineteen cases are excluded by the obstruction theorems of
Parisi--Spahiu--Skandera--Wang. Larson's earlier small-resolution
construction for $463152$, together with the framework of Agrawal and
Sotirov, gives a prior geometric route to an associated
Kazhdan--Lusztig factorization. The present note gives direct
reversal-factorization certificates in the normalization of the
conjecture and the complete $\Ssix$ classification.
\end{abstract}

\maketitle

\section{The counterexamples}

Let $H_n(q)$ be the type-$A$ Hecke algebra over $\mathbb Z[q]$ with
basis $\{T_w:w\in\Sn\}$ and multiplication
\[
T_xT_{s_i}=
\begin{cases}
T_{xs_i},&\ell(xs_i)>\ell(x),\\
qT_{xs_i}+(q-1)T_x,&\ell(xs_i)<\ell(x).
\end{cases}
\]
We use the modified signless Kazhdan--Lusztig basis
\cite{KazhdanLusztig}
\[
\Ct_w=\sum_{v\leq w}P_{v,w}(q)T_v.
\]
For $1\leq a<b\leq n$, let $s[a,b]$ be the longest element of the
standard parabolic subgroup generated by $s_a,\ldots,s_{b-1}$, and put
\[
\Cab{a}{b}:=\Ct_{s[a,b]}.
\]

For a sequence of intervals
$\mathcal A=([a_1,b_1],\ldots,[a_m,b_m])$, define
\[
I_{i,j}=[a_i,b_i]\cap[a_j,b_j]\setminus
\bigcup_{i<k<j}[a_k,b_k],
\qquad
f_{\mathcal A}(q)=\prod_{i<j}[|I_{i,j}|]_q!,
\]
where $[r]_q!=\prod_{k=1}^r(1+q+\cdots+q^{k-1})$.
Following \cite[Sec.~3]{ParisiEtAl}, a \emph{reversal factorization}
of $\Ct_w$ is an identity
\[
f_{\mathcal A}(q)\Ct_w
 =\Cab{a_1}{b_1}\cdots\Cab{a_m}{b_m}.
\]
Conjecture~5.4 of \cite{ParisiEtAl} asserts that no such factorization
exists when $w$ contains the pattern $45312$.

\begin{theorem}\label{thm:main}
Among the $26$ permutations in $\Ssix$ containing $45312$, the element
$\Ct_w$ has a reversal factorization if and only if
\[
w\in\mathcal F:=
\{463152,526413,463512,536412,562413,563142,563412\}.
\]
The seven factorizations are displayed in
Table~\ref{tab:factorizations}. In particular, Conjecture~5.4 of
\cite{ParisiEtAl} is false.
\end{theorem}

\begin{table}[t]
\centering
\small
\caption{The seven counterexamples. In every row the displayed
interval sequence satisfies $f_{\mathcal A}(q)=(1+q)^2$.}
\label{tab:factorizations}
\begin{tabular}{@{}ccl@{}}
\toprule
$w$ & positions of $45312$ & factorization in $H_6(q)$ \\
\midrule
$463152$ & $1,2,3,4,6$ &
$(1+q)^2\Ct_{463152}=\Cab{2}{4}\Cab{1}{3}\Cab{4}{6}\Cab{2}{4}$ \\
$526413$ & $1,3,4,5,6$ &
$(1+q)^2\Ct_{526413}=\Cab{3}{5}\Cab{4}{6}\Cab{1}{3}\Cab{3}{5}$ \\
$463512$ & $1,2,3,5,6$ &
$(1+q)^2\Ct_{463512}=\Cab{2}{4}\Cab{1}{3}\Cab{4}{6}\Cab{2}{4}\Cab{4}{5}$ \\
$536412$ & $1,3,4,5,6$ &
$(1+q)^2\Ct_{536412}=\Cab{2}{3}\Cab{3}{5}\Cab{1}{3}\Cab{4}{6}\Cab{3}{5}$ \\
$562413$ & $1,2,4,5,6$ &
$(1+q)^2\Ct_{562413}=\Cab{3}{5}\Cab{4}{6}\Cab{1}{3}\Cab{3}{5}\Cab{2}{3}$ \\
$563142$ & $1,2,3,4,6$ &
$(1+q)^2\Ct_{563142}=\Cab{4}{5}\Cab{2}{4}\Cab{4}{6}\Cab{1}{3}\Cab{2}{4}$ \\
$563412$ & $1,2,3,5,6$ &
$(1+q)^2\Ct_{563412}=\Cab{4}{5}\Cab{2}{4}\Cab{1}{3}\Cab{4}{6}\Cab{2}{4}\Cab{4}{5}$ \\
\bottomrule
\end{tabular}
\end{table}

\begin{proof}
Direct multiplication in the $T$-basis gives
\begin{align}
\Cab{2}{4}\Cab{1}{3}
  &=(1+q)\Ct_{431256},\label{eq:first}\\
\Ct_{431256}\Cab{4}{6}
  &=\Ct_{431652},\label{eq:second}\\
\Ct_{431652}\Cab{2}{4}
  &=(1+q)\Ct_{463152}.
\label{eq:third}
\end{align}
For the interval sequence
$([2,4],[1,3],[4,6],[2,4])$, the only $I_{i,j}$ of cardinality greater
than one are
$I_{1,2}=I_{2,4}=\{2,3\}$; hence
$f_{\mathcal A}(q)=(1+q)^2$. Equations
\eqref{eq:first}--\eqref{eq:third} therefore give the first row of
Table~\ref{tab:factorizations}.

Two further exact identities are
\[
\Ct_{463152}\Cab{4}{5}=\Ct_{463512},
\qquad
\Cab{4}{5}\Ct_{463512}=\Ct_{563412}.
\]
They give the rows for $463512$ and $563412$. The standard
anti-involution and diagram automorphism satisfy
\[
\iota(\Ct_w)=\Ct_{w^{-1}},
\qquad
\omega(\Ct_w)=\Ct_{w_0ww_0},
\qquad w_0=654321,
\]
where $\iota$ reverses products and
$\omega(\Cab{a}{b})=\Cab{7-b}{7-a}$. Applying these symmetries gives
the remaining four rows; the three resulting orbits are
\[
\{463152,526413\},\qquad
\{463512,536412,562413,563142\},\qquad
\{563412\}.
\]
A direct inspection of the interval sequences gives
$f_{\mathcal A}(q)=(1+q)^2$ in every row.

It remains to exclude the other nineteen cases. Thirteen of them are
\begin{multline*}
\mathcal N_1=\{156423,256413,453126,453162,456312,463125,465312,\\
516423,546312,564123,564132,564213,564312\}.
\end{multline*}
Every element of $\mathcal N_1$ avoids $4231$, contains $3412$, and has
$\gap(w)>1$, where $\gap$ is the $3412$-gap of
\cite[Sec.~5]{ParisiEtAl}. The gap is $2$ for the first twelve
permutations and $3$ for $564312$. Theorem~5.2 of
\cite{ParisiEtAl} therefore rules out reversal factorizations for all
of them.

For the remaining six permutations, the following pairs satisfy
$v<w$ and have the indicated Kazhdan--Lusztig polynomial:
\[
\begin{array}{ccl@{\qquad}ccl}
 w&v&P_{v,w}(q)&w&v&P_{v,w}(q)\\[2pt]
356412&213456&1+q^2&453612&123465&1+q^2\\
561423&213456&1+q^2&563124&123465&1+q^2\\
564231&134526&1+q^2&645312&152346&1+q^2.
\end{array}
\]
The coefficient of $q$ is an internal zero, so Theorem~5.3 of
\cite{ParisiEtAl} excludes a reversal factorization in every case.
The three sets above have sizes $7$, $13$, and $6$, and exhaust the
$26$ permutations in $\Ssix$ containing $45312$.

All displayed identities, pattern enumerations, gap values, and
Kazhdan--Lusztig polynomials were checked with exact arithmetic in
$\mathbb Z[q]$. The basis was generated by the standard
Kazhdan--Lusztig recurrence \cite[Chap.~5]{BjornerBrenti}, and the
products were evaluated from the defining Hecke relation; no numerical
specialization was used. A self-contained verification script,
\texttt{verify\_reversal\_factorizations.py}, using only the Python
standard library, accompanies the source.
\end{proof}

\begin{corollary}\label{cor:alln}
For every $n\geq6$, Conjecture~5.4 of \cite{ParisiEtAl} has a
counterexample in $\mathfrak S_n$.
\end{corollary}

\begin{proof}
Embed $\Ssix$ as the standard parabolic subgroup of $\mathfrak S_n$
fixing $7,\ldots,n$. Compatibility of the Kazhdan--Lusztig basis with
standard parabolic subgroups preserves the first factorization in
Table~\ref{tab:factorizations}, while
$4631527\cdots n$ still contains $45312$.
\end{proof}

\begin{remark}[Prior work and scope]\label{rem:prior}
Larson \cite[Ex.~6.11]{Larson} constructed a small
Gelfand--MacPherson resolution for $w=463152$ with resolution data
\[
I_0=\{2,3,5\},\qquad I_1=\{1,2,4,5\},\qquad
I_2=\{2,3,5\},
\]
and overlaps $J_1=J_2=\{2,5\}$. Under the correspondence developed
by Agrawal and Sotirov
\cite[Cor.~6.21 and Sec.~8]{AgrawalSotirov}, these parabolic data
yield a Kazhdan--Lusztig factorization. In the present normalization,
splitting the disconnected parabolics gives
\[
(1+q)^4\Ct_{463152}
 =\Cab{2}{4}\Cab{5}{6}\Cab{1}{3}\Cab{4}{6}
  \Cab{2}{4}\Cab{5}{6}.
\]
Since $\Cab{5}{6}$ commutes with $\Cab{1}{3}$ and $\Cab{2}{4}$, and
\[
\Cab{5}{6}\Cab{4}{6}
 =\Cab{4}{6}\Cab{5}{6}
 =(1+q)\Cab{4}{6},
\]
this contracts to the first identity of
Table~\ref{tab:factorizations}. Thus the geometric existence of the
$463152$ example is implicit in prior work.
Theorem~\ref{thm:main} records explicit reversal-factorization
certificates in the precise form of Conjecture~5.4 and gives the
complete classification of its $\Ssix$ instances. The broader
characterization problem posed in \cite[Quest.~4.5]{Skandera}
remains open.
\end{remark}

{\small
\begin{thebibliography}{99}

\bibitem{AgrawalSotirov}
R.~Agrawal and V.~Sotirov,
\emph{A generalization of Deodhar's framework for questions in
Kazhdan--Lusztig theory},
arXiv:2107.08522, 2021.

\bibitem{BjornerBrenti}
A.~Bj\"orner and F.~Brenti,
\emph{Combinatorics of Coxeter groups},
Graduate Texts in Mathematics, vol.~231, Springer, New York, 2005.

\bibitem{KazhdanLusztig}
D.~Kazhdan and G.~Lusztig,
\emph{Representations of Coxeter groups and Hecke algebras},
Invent. Math. \textbf{53} (1979), no.~2, 165--184.

\bibitem{Larson}
S.~J. Larson,
\emph{Decompositions of Schubert varieties and small resolutions},
Comm. Algebra \textbf{52} (2024), no.~5, 2101--2126.
\href{https://doi.org/10.1080/00927872.2023.2281602}
{doi:10.1080/00927872.2023.2281602}.

\bibitem{ParisiEtAl}
T.~Parisi, B.~Spahiu, M.~Skandera, and J.~Wang,
\emph{On Kazhdan--Lusztig basis elements having no reversal
factorization},
arXiv:2605.21733v1, 2026.

\bibitem{Skandera}
M.~Skandera,
\emph{On the dual canonical and Kazhdan--Lusztig bases and
$3412$-, $4231$-avoiding permutations},
J. Pure Appl. Algebra \textbf{212} (2008), no.~5, 1086--1104.
\href{https://doi.org/10.1016/j.jpaa.2007.09.007}
{doi:10.1016/j.jpaa.2007.09.007}.

\end{thebibliography}
}

\end{document}